\documentclass[10pt,a4paper]{amsart}
\usepackage{amsmath,amssymb,amsthm,graphicx}
% Retained preamble for the Stacks Project chapter {CH} (shared project preamble).
% The chapter file pinned-original-source.tex only carries \input{{preamble}} on its first line,
% so no preamble text lives inside the chapter itself.  The Stacks Project uses one single shared
% preamble for every chapter: the snapshot stages it three times (sources/stacks-chow/preamble.tex,
% sources/stacks-etale-cohomology/preamble.tex and focused-seeds/AGP/sources/stacks-snapshot-20261011/
% preamble.tex) and those copies are identical in size.  The block below is copied mechanically,
% line for line, from sources/stacks-chow/preamble.tex lines 14-18 and 154-217: the Xy-pic setup and
% the native theorem declarations together with the mathematical macros.  Nothing was added, removed
% or redefined; the wrapper class and the amsmath/amssymb/amsthm/graphicx base packages are supplied
% by the assembler.
% For commutative diagrams we use Xy-pic
% Chapter: stacks-cohomology.tex -- Cohomology of Sheaves.
\usepackage[all]{xy}

% We use 2cell for 2-commutative diagrams.
\xyoption{2cell}

% Theorem environments.
%
\theoremstyle{plain}
\newtheorem{theorem}[subsection]{Theorem}
\newtheorem{proposition}[subsection]{Proposition}
\newtheorem{lemma}[subsection]{Lemma}

\theoremstyle{definition}
\newtheorem{definition}[subsection]{Definition}
\newtheorem{example}[subsection]{Example}
\newtheorem{exercise}[subsection]{Exercise}
\newtheorem{situation}[subsection]{Situation}

\theoremstyle{remark}
\newtheorem{remark}[subsection]{Remark}
\newtheorem{remarks}[subsection]{Remarks}

\numberwithin{equation}{subsection}

% Macros
%
\def\lim{\mathop{\mathrm{lim}}\nolimits}
\def\colim{\mathop{\mathrm{colim}}\nolimits}
\def\Spec{\mathop{\mathrm{Spec}}}
\def\Hom{\mathop{\mathrm{Hom}}\nolimits}
\def\Ext{\mathop{\mathrm{Ext}}\nolimits}
\def\SheafHom{\mathop{\mathcal{H}\!\mathit{om}}\nolimits}
\def\SheafExt{\mathop{\mathcal{E}\!\mathit{xt}}\nolimits}
\def\Sch{\mathit{Sch}}
\def\Mor{\mathop{\mathrm{Mor}}\nolimits}
\def\Ob{\mathop{\mathrm{Ob}}\nolimits}
\def\Sh{\mathop{\mathit{Sh}}\nolimits}
\def\NL{\mathop{N\!L}\nolimits}
\def\CH{\mathop{\mathrm{CH}}\nolimits}
\def\proetale{{pro\text{-}\acute{e}tale}}
\def\etale{{\acute{e}tale}}
\def\QCoh{\mathit{QCoh}}
\def\Ker{\mathop{\mathrm{Ker}}}
\def\Im{\mathop{\mathrm{Im}}}
\def\Coker{\mathop{\mathrm{Coker}}}
\def\Coim{\mathop{\mathrm{Coim}}}

% Boxtimes
%
\DeclareMathSymbol{\boxtimes}{\mathbin}{AMSa}{"02}

%
% Macros for moduli stacks/spaces
%
\def\QCohstack{\mathcal{QC}\!\mathit{oh}}
\def\Cohstack{\mathcal{C}\!\mathit{oh}}
\def\Spacesstack{\mathcal{S}\!\mathit{paces}}
\def\Quotfunctor{\mathrm{Quot}}
\def\Hilbfunctor{\mathrm{Hilb}}
\def\Curvesstack{\mathcal{C}\!\mathit{urves}}
\def\Polarizedstack{\mathcal{P}\!\mathit{olarized}}
\def\Complexesstack{\mathcal{C}\!\mathit{omplexes}}
% \Pic is the operator that assigns to X its picard group, usage \Pic(X)
% \Picardstack_{X/B} denotes the Picard stack of X over B
% \Picardfunctor_{X/B} denotes the Picard functor of X over B
\def\Pic{\mathop{\mathrm{Pic}}\nolimits}
\def\Picardstack{\mathcal{P}\!\mathit{ic}}
\def\Picardfunctor{\mathrm{Pic}}
\def\Deformationcategory{\mathcal{D}\!\mathit{ef}}
\usepackage{cleveref}
\crefname{originallemma}{original lemma}{original lemmas}
\Crefname{originallemma}{Original lemma}{Original lemmas}
\makeatletter
\newlabel{derived-lemma-compose-derived-functors}{{015M}{source}{}{sourcepointer-1}{}}
\newlabel{derived-lemma-compose-derived-functors@cref}{{[originallemma][0][]015M}{[1][0][]source}{}{}{}}
\makeatother
\crefname{originallemma}{original lemma}{original lemmas}
\Crefname{originallemma}{Original lemma}{Original lemmas}
\makeatletter
\newlabel{derived-lemma-right-derived-exact-functor}{{015F}{source}{}{sourcepointer-2}{}}
\newlabel{derived-lemma-right-derived-exact-functor@cref}{{[originallemma][0][]015F}{[1][0][]source}{}{}{}}
\makeatother
\crefname{originallemma}{original lemma}{original lemmas}
\Crefname{originallemma}{Original lemma}{Original lemmas}
\makeatletter
\newlabel{lemma-pushforward-injective}{{01EX}{source}{}{sourcepointer-3}{}}
\newlabel{lemma-pushforward-injective@cref}{{[originallemma][0][]01EX}{[1][0][]source}{}{}{}}
\makeatother
\crefname{originalremark}{original remark}{original remarks}
\Crefname{originalremark}{Original remark}{Original remarks}
\makeatletter
\newlabel{remark-elucidate-lemma}{{01F0}{source}{}{sourcepointer-4}{}}
\newlabel{remark-elucidate-lemma@cref}{{[originalremark][0][]01F0}{[1][0][]source}{}{}{}}
\makeatother
\title{Commutativity of derived global sections with the derived pushforward and the restriction of scalars}
\author{The Stacks Project Authors}
\date{stacks-project:cohomology.tex:local-snapshot-20261011; GNU FDL 1.2 or later}
\begin{document}
\maketitle
\noindent\textit{Editorial note.} All delivered mathematical text is the author's own, copied from the pinned Stacks Project chapter file cohomology.tex; no statement or proof step was rewritten, reordered or removed.\par\medskip
\noindent\textit{Editorial note.} The statement ends by pointing the reader to a remark for more explanation; the remark is carried as a navigation stub and is not part of the proof, which is complete as written.\par\medskip
\noindent\textit{Editorial note.} The counted spans carry 4 navigation stubs and no bibliography citation.\par\medskip
\noindent\textit{Editorial note.} The preamble is the chapter's retained minimal preamble (sources/stacks-cohomology/retained-preamble.tex lines 1-81), a verbatim excerpt of the native preamble.\par\medskip
\noindent\textit{Editorial note.} The label lemma-before-Leray is defined only in cohomology.tex, so no cross-chapter restatement had to be weighed for this card.\par\medskip
\noindent\textit{Allowed standard prerequisites (editorial).}\begin{itemize}
\item Right derived functor of an exact functor is the functor itself. Derived Categories, Lemma derived-lemma-right-derived-exact-functor, invoked at proof lines 2125-2128. Supplies that the derived restriction functor is computed by applying the functor.
\item Composition of derived functors and the acyclicity condition. Derived Categories, Lemma derived-lemma-compose-derived-functors, invoked at proof lines 2137-2139. Supplies the commutativity once the hypotheses are verified for both compositions.
\item Pushforward of an injective module is acyclic for global sections. Same chapter, Lemma lemma-pushforward-injective, invoked at proof lines 2141-2143. Supplies the acyclicity hypothesis for the second composition.
\end{itemize}\medskip
\ifdefined\hypertarget\hypertarget{sourcepointer-1}{}\fi
\noindent\textit{Source reference.} Original lemma 015M: sources/stacks-chow/tags.txt, source lines 1513--1513. Pointer outside the counted span, invoked once for context; the load-bearing argument of the counted statement is the author's proof reproduced in full in the proof span.\par\medskip
\ifdefined\hypertarget\hypertarget{sourcepointer-2}{}\fi
\noindent\textit{Source reference.} Original lemma 015F: sources/stacks-chow/tags.txt, source lines 1506--1506. Pointer outside the counted span, invoked once for context; the load-bearing argument of the counted statement is the author's proof reproduced in full in the proof span.\par\medskip
\ifdefined\hypertarget\hypertarget{sourcepointer-3}{}\fi
\noindent\textit{Source reference.} Original lemma 01EX: sources/stacks-cohomology/pinned-original-source.tex, source lines 1779--1779. Pointer outside the counted span, invoked once for context; the load-bearing argument of the counted statement is the author's proof reproduced in full in the proof span.\par\medskip
\ifdefined\hypertarget\hypertarget{sourcepointer-4}{}\fi
\noindent\textit{Source reference.} Original remark 01F0: sources/stacks-cohomology/pinned-original-source.tex, source lines 2147--2147. Pointer outside the counted span, invoked once for context; the load-bearing argument of the counted statement is the author's proof reproduced in full in the proof span.\par\medskip
\begin{lemma}
\label{lemma-before-Leray}
Let $f : X \to Y$ be a morphism of ringed spaces.
There is a commutative diagram
$$
\xymatrix{
D^{+}(X) \ar[rr]_-{R\Gamma(X, -)} \ar[d]_{Rf_*} & &
D^{+}(\mathcal{O}_X(X)) \ar[d]^{\text{restriction}} \\
D^{+}(Y) \ar[rr]^-{R\Gamma(Y, -)} & &
D^{+}(\mathcal{O}_Y(Y))
}
$$
More generally for any $V \subset Y$ open and $U = f^{-1}(V)$ there
is a commutative diagram
$$
\xymatrix{
D^{+}(X) \ar[rr]_-{R\Gamma(U, -)} \ar[d]_{Rf_*} & &
D^{+}(\mathcal{O}_X(U)) \ar[d]^{\text{restriction}} \\
D^{+}(Y) \ar[rr]^-{R\Gamma(V, -)} & &
D^{+}(\mathcal{O}_Y(V))
}
$$
See also Remark \ref{remark-elucidate-lemma} for more explanation.
\end{lemma}

\begin{proof}
Let
$\Gamma_{res} : \textit{Mod}(\mathcal{O}_X) \to \text{Mod}_{\mathcal{O}_Y(Y)}$
be the functor which associates to an $\mathcal{O}_X$-module $\mathcal{F}$
the global sections of $\mathcal{F}$ viewed as an $\mathcal{O}_Y(Y)$-module
via the map $f^\sharp : \mathcal{O}_Y(Y) \to \mathcal{O}_X(X)$. Let
$restriction : \text{Mod}_{\mathcal{O}_X(X)} \to \text{Mod}_{\mathcal{O}_Y(Y)}$
be the restriction functor induced by
$f^\sharp : \mathcal{O}_Y(Y) \to \mathcal{O}_X(X)$. Note that $restriction$
is exact so that
its right derived functor is computed by simply applying the restriction
functor, see
Derived Categories, Lemma \ref{derived-lemma-right-derived-exact-functor}.
It is clear that
$$
\Gamma_{res}
=
restriction \circ \Gamma(X, -)
=
\Gamma(Y, -) \circ f_*
$$
We claim that
Derived Categories, Lemma \ref{derived-lemma-compose-derived-functors}
applies to both compositions. For the first this is clear by our remarks
above. For the second, it follows from
Lemma \ref{lemma-pushforward-injective} which implies that
injective $\mathcal{O}_X$-modules are mapped to $\Gamma(Y, -)$-acyclic
sheaves on $Y$.
\end{proof}

\end{document}
